Was sind die vorvertraglichen Beratungspflichten des Versichers?

Befragungs, Beratungs- einschließlich Begründungspflichten sowie Dokumentationspflichten. 
Beratungsadressat ist der Versicherungsnehmer; 
dem Versicherungsvermittler bestehen eigene entsprechende Pflichten. Für § 61 VVG sind ggf. weitere Pflichten aus dem Maklervertrag.

Befragungspflichten:
Die vorvertragliche Fragepflicht hat sich nach den Wünschen und Bedürfnissen des Versicherungsnehmers auszurichten. 
Auf dieser Beratungsgrundlage hat der Versicherer den Versicherungsnehmer in angemessenem Verhältnis von Beratungsaufwand und der vom 
Versicherungsnehmer zu zahlenden Prämie zu beraten. 

Dokumentationspflichten:
Die Gründe für jeden zu einer bestimmten Versicherung erteilten Rat sind anzugeben. 
Bei der diesbezüglichen Dokumentierung hat der Versicherer die Komplexität des angebotenen Versicherungsvertrags zu berücksichtigen (§ 6 I1 VVG). 
Vor Abschluss des Versicherungsvertrags sind der erteilte Rat und die Gründe klar und verständlich in Papierform oder auf einem Versicherers mit 
personalisiertem Zugang verfügbar zu machen (VVG §6a I).

Sie haben somit eine Doppelfunktion:
• Erhöhen der Beratungsqualität
• Bewältigung von Beweisproblemen bzgl Fehlberatung
und sind nach Formalia grundsätzlich in Textform.

Der Verzicht auf Beratung ist nur unter eindeutigem Hinweis auf mögliche Benachteiligung bei Schadensersatzansprüchen gegenüber dem Versicherer
möglich.


Was sind mögliche Konsequenzen bei Verletzung der Beratungspflicht?

Eine Pflichtverletzung liegt u.a. in der unvollständigen, verspäteten oder falschen Beratung des Versicherungsnehmers. 
Der Versicherer hat sich gemäß § 278 S. 1 BGB das Verhalten von Angestellten und Versicherungsvertretern als seinen Erfüllungsgehilfen zurechnen zu lassen. 
Kein Erfüllungsgehilfe ist der Versicherungsmakler, da er im Lager des Versicherungsnehmers steht (vgl. auch § 6 VI VVG). 

§280 I BGB nachgebildet (auch Beweislast in S. 2); § 254 BGB gilt
• Entsprechend auch Zurechnung und Haftung für selbst. und unselbst. VersVertreter
(Erfüllungsgehilfe § 278 BGB), halbzwingend gem. § 18 VVG!
• Vertragsaufhebung, wenn Schaden = nachteiliger Vertrag und Vers-Fall noch nicht
eingetreten (Ersatz der Prämien und der Aufwendungen für den Vertragsschluss)
• Vertragsanpassung, wenn VR bedarfsgerechten Vertrag anbietet (Deckungsschutz
unter Anrechnung der Differenzprämie)
• Wäre Risiko bei ordnungsgemäßer Beratung einbezogen worden:
„Erfüllungshaftung“ (i.E. str.)

Beweiserleichterungen
• für den VN bei Verletzung von Dokumentationspflichten (z.B. hinsichtlich gestellter
Fragen) etwa Nichtvorlage Beratungsprotokoll


Der Versicherungsnehmer bekommt den Vertrauensschaden ersetzt, der bis zur „Quasi-Deckung" führen kann. 

Was ist die Beratungspflicht während der Vertragslaufzeit?

Nach § 6 IV VVG sind Anlässe
• Einführung neuer AVB/Tarife, aber nur im Rahmen von Verhandlungen (str.)
• Eintritt des Versicherungsfalls (aber nicht: Obliegenheiten, str.)
• Unwirksamkeit einer Kündigung
• Vertragswechsel in der PKV oder BUV
• Bei konkreten Wünschen des VN nur, wenn Wunsch erkennbar von

Fehlvorstellung getragen
• Dauer
• bis zum Vertragsende
• Keine Dokumentationspflicht
• Dokumentation aber zu Beweiszwecken sinnvoll



Erläutern Sie die Gestaltung der Rückwärtsversicherung und vorläufiger Deckung

Im Gegensatz zum klassischen Versicherungsbeginn (e.g. wenn Dauer nach Tagen, Wochen, Monaten bestimmt: 
beginnt mit Beginn des Tages des Vertragsschlusses = 0:00 Uhr des Vertragsschlusstages, sogn. „Mitternachtsregel“,
oder mit Eingang erster Prämienzahlung unter ausdrücklichem Hinweis, siehe § 10 VVG)

Es ist zu differenzieren in 
Rückwärtsversicherung
- setzt Zustandekommen des Hauptvertrages voraus
Vorläufige Deckung 
- besteht auch dann, wenn es nicht zu einem Hauptvertrag kommt (z.B. Versicherer lehnt Antrag nach Risikoprüfung ab)

Die vorläufige Deckung ist typisch für die Kfz-Haftpflichtversicherung.
Rechtliche Grundlage dafür ist die erleichterte Form für den Abschluss eines vorläufigen Deckungsvertrages nach §49 I VVG.
Durch ein konkludenten Verzicht auf die rechtzeitige  Übermittlung der Vertragsbestimmungen, sprich statdessen erst zusammen mit Versicherungsschein,
nimmt der Versicherungsnehmer die üblicherweise verwendeten Bedingungen des Versicherers an.

Nach §51 I VVG darf Versicherer dies abhängig von einer ersten Prämienzahlung/Einmalzahlung machen. Dies ist in der Praxis unwesentlich.
Dem gegenüber steht der rückwirkende Wegfall des Versicherungsschutzes, wenn der Versicherungsnehmer scheitert die Prämie nach Erhalt des Versicherungsscheines
innerhalb von zwei Wochen (wenn dies vertretbar ist) zu zahlen. Beweispflicht liegt hier beim Versicherer.



Was ist die IDD und ihr Anwendungsbereich?

Insurance Distribution Directive / Versicherungsvertriebsrichtlinie

Was ist Versicherungsvermittlung nach IDD?
Art. 2 Abs. 1 Nr. 1 IDD: 
Beratung, Vorschlagen oder Durchführen anderer Vorbereitungshandlungen zum Abschließen von Versicherungsverträgen
Abschließen von Versicherungsverträgen
Mitwirken bei der Verwaltung und Erfüllung von Versicherungsverträgen, insbesondere im Schadensfall

Keine Versicherungsvermittlung:
Berufsmäßige Verwaltung der Ansprüche eines VU
Schadenregulierung
Sachverständigenbegutachtung von Schäden
„Gruppenorganisator“ (Gruppen-VN) von Gruppenversicherungsverträgen 
Beachte EuGH Urteil vom 29.09.2022 (Az.: C-633/20) 


Was sind die Rechtsformen von Versicherungsvermittlung?

Anwendbar auf „jede natürliche oder juristische Person, die in einem Mitgliedstaat niedergelassen ist 
oder sich dort niederlassen will, um den Vertrieb von Versicherungs- und Rückversicherungsprodukten aufzunehmen oder auszuüben“
(Inbesondere schließt dies Tippgeber aus)

Versicherungsvertreter
- Ausschließlichkeit / hauptberuflich
- Ausschließlichkeit / nebenberuflich
- Mehrfachvertreter / hauptberuflich
- Mehrfachvertreter / nebenberuflich
Versicherungsmakler
VersVermittler mit Erlaubnisbefreiung (produktakzessorischer Vermittler)
Gebundener Versicherungsvertreter
Annexvermittler / Bagatellvermittler
Angestellte des Versicherers / Handelsgehilfe

VersVermittler (Hauptberuflich)
Wer gewerbsmäßig (gegen Vergütung!) den Abschluss von Versicherungs- oder Rückversicherungsverträgen vermitteln will (Versicherungsvermittler), 
bedarf nach Maßgabe der folgenden Bestimmungen der Erlaubnis der zuständigen Industrie- und Handelskammer. 
Versicherungsvermittler ist, wer 
als Versicherungsvertreter eines oder mehrerer Versicherungsunternehmen oder eines Versicherungsvertreters damit betraut ist, 
Versicherungsverträge zu vermitteln oder abzuschließen
(Dies schließt aus dass sie ein Versicherungsunternehmen sind)
Rückversicherungsvermittler fallen nicht unter diese Richtlinien

Versicherungsmakler (erlaubnispflichtig)
Als Versicherungsmakler für den Auftraggeber die Vermittlung oder den Abschluss von Versicherungsverträgen übernimmt, 
ohne von einem Versicherungsunternehmen oder einem Versicherungsvertreter damit betraut zu sein.

Gebundener Vermittler/ Ausschließlichkeitsvertreter(erlaubnisfrei)
Tätigkeit als Versicherungsvermittler ausschließlich im Auftrag eines oder, wenn die Versicherungsprodukte nicht in Konkurrenz stehen, 
mehrerer Versicherungsunternehmen ausübt, die für ihn sämtliche Haftung übernehmen.


Produktakzessorischer Vermittler (erlaubnispflichtig / auf Antrag erlaubnisfrei)
Gewerbetreibender, der die Versicherung als Ergänzung der im Rahmen seiner Haupttätigkeit gelieferten Waren oder Dienstleistungen vermittelt, 
wenn
- Person vertreibt lediglich bestimmte Versicherungsprodukte, die eine Ergänzung zur Lieferung einer Ware bzw. zur Erbringung einer Dienstleistung darstellen
- keine Lebensversicherungs- oder Haftpflichtrisiken (außer es ist wirklich wesentlich für Hauptprodukt)
- er zuverlässig sowie angemessen qualifiziert ist und nicht in ungeordneten Vermögensverhältnissen lebt.
- registrierungspflichtig

Produktakzessorischer Vermittler / Bagatellvermittler (erlaubnisfrei)
Versicherungsvermittler in Nebentätigkeit, wenn sämtliche nachstehenden Bedingungen erfüllt sind:
Versicherung ist ergänzende Leistung zur Lieferung einer Ware bzw. zur Erbringung einer Dienstleistung Versicherung deckt Folgendes ab:
Defekts, Verlust oder Beschädigung der Ware oder der Nichtinanspruchnahme der Dienstleistung
Beschädigung oder Verlust von Gepäck und andere Risiken im Zusammenhang mit einer bei dem betreffenden Anbieter gebuchten Reise
Jahresprämie < 600 EUR
Prämie bei 3 Monate Laufzeit < 200 EUR


Wie unterscheiden sich Vermittler und Vertreter in Verantwortung und Haftbarkeit? 

Vertreter sind in §59 II VVG definiert: Von VR mit gewerbsmäßiger Versicherungsvermittlung betraut
(ähnlich zu 84 ff. HGB) und damit nicht Angestellte des VR.
Die Relation ist meist durch einen Agenturvertrag (= Dienstvertrag) geregelt
Provisionsanspruch gegen VR nach § 92 III,IV HGB, wenn VN Prämie zahlt
Nach § 89b HGB Ausgleichsanspruch bei Ende des Agenturvertrags.
Vertreter haben Eigenhaftung aus §63 VVG bei Verletzung von Pflichten wie der Beratung/Dokumentation.

Ihre Vertretungsmacht beläuft sich abhängig von Vermittlungs- und Abschlussvertretern auf:
Gesetzliche Vollmacht nach § 69 VVG (Abschluss von Vertrag, Übermittlung Versicherungsschein, Inkassovollmacht (Prämienannahme))

Der Versicherungsmakler nach § 59 III VVG ist nicht von VR betraut sondern leistet für den VN
via Maklervertrag. Er hat keine gesetzliche Vollmacht ist aber vetraglich meist mit dieser betraut.
Pflichten: Umfassende Betreuung der Versicherungsinteressen des VN,
umfassende Marktanalyse
Haftung: Pflichtverletzung §§ 60,61,63

Er wird abweichend von § 99 HGB nicht von beiden Parteien vergütet, sondern nur von VR was meist als Margin
in die Prämie verrechnet wird (Bruttopolice). Seltener gibt es auch eine Courtage von seiten des VN (dann Nettopolice).

Auffällig ist das der Makler häufig als Wissenvertreter (§§ 164, 166 BGB) für den VN einsteht und dem VR eine Wissenserklärung abgibt
(Verwaltung der Anzeige und Aufklärungsobliegenheiten).


Wie sind Dritte im Versicherungsvertrag reregelt?

Es ist zu unterscheiden zwischen Versicherung für fremde Rechnung und Bezugsberechtigten

Beispiel für fremde Rechnung: Leasingnehmer (VN) schließt Kfz-Haftpfl. ab, dann ist Leasinggeber Versicherter.
Grundlage für fremde Rechnung ist §§43 VVG
Die Rechte aus dem Vertrag stehen dem Versicherten zu. Hingegen kann nur der VN über Ansprüche verfügen und sie gerichtlich geltend machen.
Der Versicherungsschein wird dem Versicherungsnehmer ausgehändigt, dieser kann der Erlaubnis zur Weitergabe an Versicherten erteilen
damit dieser rechtliche Schritte ebenfalls einleiten kann. Obliegenheiten und Risikoausschlüsse treffen beide.

Bezugsberechtigte sind komplexer und häufig im Bereich der Lebensversicherungen anzutreffen.


Welche Konsequenzen und Regelungen gelten für die Gefahrerhöhung?

Grundsätzlich darf nach §23 VVG der VN nach Vertragserklärung ohne Einwilligung des VR keine Gefahrenerhöhung 
vornehmen. Die Pflichten ergeben sich im Detail auch aus §23 VVG.
Der Versicherer nimmt die nach §19 erforderliche Risikoprüfung vor welche die Geschäftsgrundlage darstellt. 
Erkennt der VN nachträglich dass er eine Gefahrerhöhung vorgenommen hat, so ist diese nach §23 II VVG unverzüglich anzuzeigen
(Anzeigeobliegenheit), ebenfalls wenn diese durch Dritte ensteht (§23 III).
Unter bestimmten Voraussetzung kann der VR nach 
- §24 I den Vertrag kündigen
- §25 I eine Vertragsanpassung vornehmen. 
für Versicherungsfälle nach Erhöhung kann er leistungsfrei bleiben wenn § 26 greift.



Definieren Sie Obliegenheiten und differenzieren Sie von ähnlichen Konzepten 

Der Versicherungsnehmer hat vor und nach Vertragsabschluss Obliegenheiten zu erfüllen. 
Obliegenheiten sind Verhaltensnormen, die 
- den Vertragszweck wahren sollen 
- bei Nichtbeachtung unter Voraussetzungen zur Leistungsfreiheit führen

Im Versicherungsrecht ist zu unterscheiden zwischen folgenden Obliegenheiten 
- vertraglichen (§ 28 VVG)
- gesetzlichen. 
Beispiele für Gesetzlichen
- vorvertragliche Anzeigepflicht (§§ 16 ff. VVG)
- Gefahrerhöhung (§§ 23 ff.)
- Rettungspflicht nach § 62 VVG.

Mit Ausnahme der Anzeige- und Auskunftsobliegenheit (§§ 30, 31) gilt § 28 (insbesondere seine Konsequenzen) nicht für Gesetzliche.

Die Berufung des Versicherers auf Leistungsfreiheit oder Leistungskürzung wegen einer Obliegenheitsverletzung 
ist nicht von Amts wegen zu beachten, sondern muss ausdrücklich spätestens in der 1. Instanz (selten kann 2. nachgeholt werden) 
eines Prozesses erfolgen. 
Es steht also zur Disposition des Versicherers, ob er sich auf die Obliegenheitsverletzung beruft. 
Teilweise ist die Rechtsprechung der Meinung, dass eine Geltendmachung auch noch in der 2. Instanz nachgeholt werden kann.

Sie sind abzugrenzen von
- Rechtspflichten (§ 241 BGB) denn sie sind nicht erzwingbar (= einklagbar), sie führen nicht zu Schadenersatzansprüchen
- Risikoausschlüssen (§ 103 VVG), d.h. "eh nicht versichert"
- Verhüllte Obliegenheiten: Verhaltensanforderungen an VN können nicht durch objektivierende Formulierungen zu Risikoausschlüssen gemacht werden.
Klauseln mit verhüllten Obliegenheiten welche einen harten Risikoausschluss darstellen sind eng auszulegen und können ggf.
nach §305c I BGB als überraschend oder nach §307 BGB unwirksam ausgelegt werden.


Ordnen Sie Obliegenheiten nach Zeit an

(1) Vor Abschluss des VersV zu erfüllen („Anzeigepflicht“ § 19 VVG)
(2) Nach Abschluss des VersV aber vor Eintritt des VersFalls zu
erfüllen (§§ 23ff „Gefahrerhöhung“, § 28 I VVG vertragl Obliegenheit)
• (3) Nach Eintritt des VersFalls zu erfüllen (§ 30 „Anzeige“, § 31
„Auskunft“; § 82 „Schadensminderung“, § 28 II, III VVG vertragl
Obliegenheit)


Was ist bei der Anzeigepflicht vor Abschluss und danach des VV zu beachten?

Die Anzeigepflicht dient zum Ausgleich des Informationsdefizit des VR, entsprechend hat dieser eine Prüfungspflicht
(in Textform (§ 126b BGB)).

WICHTIG!
Anzeigepflichtig sind nur Umstände, die gefahrerheblich sind und nach denen der Versicherungsnehmer gefragt wurde

Streitig ist, ist die Notwendigkeit von sehr allg. Fragen (geradezu offensichtlich) und sind im Einzelfall zu klären.
Gefahrerheblich sind Umstände, die den Entschluss des Versicherers beeinflussen können, 
den Versicherungsvertrag überhaupt oder so wie beantragt abzuschließen. 
Das sind alle objektiven und subjektiven Umstände, die den Eintritt des versicherten Risikos wahrscheinlich machen 
und den Umfang der geschuldeten Leistung erhöhen könnten. 
Die Fragen nach den gefahrerheblichen Umständen stellen keine AGB dar. 
Der Inhalt einer Frage ist durch Auslegung zu ermitteln (Grundsatz der engen Frageauslegung).

Die Nichtbeantwortung stellt keine Verneinung bzw. Falschbeantwortung dar. Vielmehr trifft den VR eine Nachfragepflicht.
Bei widersprüchlichen Angaben des VN greift diese ebenfalls, außer Sie ist zurechnbar (e.g. Fragebogen leer abgeben).

Nach Vertragsschluss ist insbesondere die Anzeigepflicht §30, Auskunftspflicht §31 relevant, sie sind "lex imperfecta"
da seine keine Konsequenzen enthalten.


Skizieren sie die Gestaltungsrechte für alle Obliegenheitsverletzungen.

Leistungsfreiheit des VR, § 28 II VVG
Erfasst alle vertragl. Obliegenheit vor und nach Eintritt des
Versicherungsfalls , sofern Parteien bei Verletzung Leistungsfreiheit
vereinbart haben. VR muss VN über Folgen von Verletzungen nach Eintritt
aber belehren, um leistungsfrei zu sein, § 28 IV VVG

Verschulden des VN bestimmt Umfang der Leistungsfreiheit:
Es ist zu unterscheiden zwischen
• Vorsatz § 28 II 1 VVG: Vollständig leistungsfrei
• Grobe Fahrlässigkeit § 28 II 2 VVG: Quotenprinzip, d.h. Leistungskürzung
entsprechend Schwere des Verschuldens (vor der VVG-Reform noch „Alles-oder-
nichts“-Prinzip)
• Einfache Fahrlässigkeit: Volle Leistungspflicht des VR

Beweislast stark differenziert nach § 28 II 2:
• Vorsatz: VR, aber sekundäre Darlegungslast des VN
• „Ob“ der groben Fahrlässigkeit: VN
• “Wie“ der groben Fahrlässigkeit, d.h. Schwere für Quotelung: VR

Durchführung der Quotelung bei grober Fahrlässigkeit:
• TdL: Mittelwertmodell: Ausgangspunkt 50\% Kürzung
• Drei-Bereichs-Modell: Nähe zur (1) einfachen Fahrlässigkeit ? (2) zum
„mittleren Bereich“? (3) zum bedingten Vorsatz ?
• BGH: In besonders schwerwiegenden Fäll auch Kürzung auf Null

Jedoch kann eine Obliegenheitsverletzung gleichgültig werden durch
einen Kausalitätsgegenbeweis, § 28 III 1 VVG:
Wenn VN vorsätzlich oder grob fahrlässig handelt kann er sich
entlasten, wenn er nachweist, dass Obliegenheitsverletzung weder
für Eintritt/Feststellung des Versicherungsfalls noch für
Feststellung/Umfang der Leistungspflicht des VR ursächlich ist
Beweislast trifft VN, dass VR keinen konkreten wirtschaftlichen
Nachteil hat.

INSERT IMAGE LECTURE 13 SLIDE 4
